\documentclass[10pt,a4paper]{article}
\usepackage[T1]{fontenc}
\usepackage[utf8]{inputenc}
\usepackage{lmodern,microtype}
\usepackage[margin=18mm]{geometry}
\usepackage{booktabs,tabularx,array,xcolor,hyperref}
\definecolor{navy}{HTML}{17365D}
\hypersetup{colorlinks=true,linkcolor=navy,urlcolor=navy}
\setlength{\parindent}{0pt}
\setlength{\parskip}{0.6\baselineskip}
\newcolumntype{Y}{>{\raggedright\arraybackslash}X}
\title{\textbf{XTSF: Synthetic Forecasting Explanations with XPC}\\Executive Summary}
\author{Thesis project}
\date{Results available through 7 August 2026}
\begin{document}
\maketitle

\section*{Decision snapshot}
Interventional Shapley converges toward the known additive structural
contribution. Conditional Shapley also converges toward its own exact estimand,
but retains a large structural mismatch because dependent features redistribute
credit. Coalitional grouping gives the best current accuracy/cost compromise;
the simplified game is fastest but changes the estimand materially.

\section{Convergence}
Increasing coalition or masking budgets reduces Monte Carlo error for both
methods. At 32 sampled coalitions, interventional MSE is 0.103 against both its
exact estimand and structural truth. Conditional MSE is 0.172 against its exact
estimand but 7.31 against structural truth. With 32 masking samples, the
corresponding values are 0.088 for interventional, and 0.333 versus 7.53 for
conditional.

The persistent conditional structural error is therefore not primarily a Monte
Carlo failure. It is a difference between the conditional attribution estimand
and the chosen signed structural decomposition.

\section{Aggregation experiment}
\begin{center}
\small
\begin{tabularx}{\linewidth}{@{}llrrrr@{}}
\toprule
Method & Mode & Structural MSE & Estimand MSE & ms/point & Rows \\
\midrule
Interventional & coalitional & 0.171 & 0.171 & 4.14 & 1024 \\
Interventional & default & 0.512 & 0.512 & 12.40 & 3072 \\
Interventional & simplified & 1.594 & 1.594 & 0.18 & 64 \\
Conditional & coalitional & 10.72 & 0.573 & 22.31 & 1024 \\
Conditional & default & 12.48 & 1.362 & 48.15 & 3072 \\
Conditional & simplified & 11.57 & 1.267 & 0.79 & 64 \\
\bottomrule
\end{tabularx}
\end{center}
Coalitional grouping improves accuracy over post-hoc default aggregation while
using one third as many model rows. The two-player simplified mode is orders of
magnitude faster but has noticeably larger estimand error and a different
interaction allocation.

\section{Supported conclusions}
\begin{itemize}
\item Monte Carlo convergence must be reported separately from estimand
      validity.
\item Exact conditional explanations are not automatically structural
      explanations under feature dependence.
\item Grouping before explanation is preferable to post-hoc aggregation in the
      current synthetic benchmark.
\item Simplified games are useful speed controls, not drop-in equivalents.
\end{itemize}

\section{Limitations and next evidence}
The evidence is synthetic and model-specific. Fixed output filenames and a
cached panel require disciplined regeneration. The next study should vary
dependence strength, background size, model class, and grouping structure while
retaining exact-estimand and structural-error diagnostics.
\end{document}
